% ===========================================================================
%  04 — 方法论
%
%  事实来源：paper_context/repo_inventory.md §2--8
%  声明：     C-TELEM-01, C-RAG-01, C-VLM-01, C-VLM-02, C-VLM-03, C-VLM-03a,
%            C-VLM-03b, C-VLM-07, C-DATA-01, C-TRAIN-01, C-PROC-01, C-PROC-02
% ===========================================================================
\chapter{方法论}\label{ch:methodology}

本章描述支撑已实现 SimTutor 系统的方法论框架。第\ref{sec:method-telemetry}节涵盖将原始 DCS-BIOS 数据转换为规范化运行时状态的遥测处理流水线。第\ref{sec:method-knowledge}节描述知识溯源子系统，包括 BM25 检索和来源策略执行。第\ref{sec:method-vlm-pipeline}节介绍从截图捕获到 LoRA 微调再到基准测试评估的七阶段视觉事实数据流水线。第\ref{sec:method-datasets}节详述六个精选数据集、从 8-fact v1 到 13-fact v2 的本体演进以及训练方案构建。第\ref{sec:method-lora-config}节说明 LoRA 微调配置，包括两个骨干模型的确切超参数。第\ref{sec:method-benchmark}节定义基准测试协议、度量和留出策略。本章为第\ref{ch:system-design}章呈现的系统设计和第\ref{ch:evaluation}章报告的技术结果提供了方法论基础。

% ===========================================================================
\section{程序运行时与遥测数据基础}\label{sec:method-telemetry}

% 证据：adapters/telemetry_pipeline.py, core/gating.py, core/vars.py
% 声明：C-TELEM-01, C-PROC-01, C-PROC-02

程序运行时是 SimTutor 跟踪学习者 F/A-18C 冷启动程序进度并生成基于证据、状态感知教学输出的核心机制。它摄入原始仿真器遥测数据，将其规范化为稳定状态变量，评估门控规则以确定步骤合格性，并将结果状态表示提供给帮助循环上下文组装器。本节描述遥测处理链的每个阶段以及门控规则评估框架。

\subsection{DCS-BIOS 原始帧摄入}

% 证据：adapters/dcs_bios/receiver.py, simtutor/schemas/v2/dcs_bios_frame.json

遥测子系统始于原始 DCS-BIOS 数据摄入。DCS-BIOS 接收器在 UDP 套接字上监听仿真器发送的二进制 DCS-BIOS 导出帧。每帧被解码为结构化 JSON 对象，附带挂钟时间戳和单调递增的序列号。解码后的帧被写入 JSONL 文件用于存档和回放目的，并转发到遥测增强流水线中。
% 证据：adapters/dcs_bios/receiver.py, simtutor/schemas/v2/telemetry_frame.json

DCS-BIOS 协议涵盖广泛的座舱控件和指示器——开关位置、旋钮状态、警告灯以及显示相关状态标志——但原始帧是嘈杂的。开关回弹、传感器漂移、快速波动的仪表值以及 UDP 乱序传递会产生虚假变化，如果直接传递给教学运行时，将产生虚假的状态变化事件并触发不必要的帮助循环。后续流水线阶段系统地解决了这些问题。

\subsection{变量解析与源丢失跟踪}

% 证据：core/vars.py, packs/fa18c_startup/telemetry_map.yaml

遥测映射（在程序包的 \texttt{telemetry\_map.yaml} 中指定）定义了原始 DCS-BIOS 字段如何投影到规范化状态变量上。变量解析器消费增强后的遥测观察，并生成一个扁平的键值对字典，其中每个键对应程序包中的一个命名变量（如 \texttt{apu\_ready}、\texttt{battery\_switch\_on}）。
% 证据：core/vars.py, packs/fa18c_startup/telemetry_map.yaml

至关重要的是，解析器在 \texttt{vars\_source\_missing} 集合中跟踪不可用或未知的来源。这一机制使下游组件能够区分"值为零"和"值未知"，这对正确的门控规则行为至关重要。缺失的变量永远不会静默地评估为成功条件。
% 证据：core/vars.py（vars_source_missing 跟踪）

\subsection{门控规则评估}

% 证据：core/gating.py, packs/fa18c_startup/pack.yaml

程序引擎为每个步骤评估两种类型的门控：前提门控（在步骤变为活跃之前必须满足）和完成门控（在步骤被标记为完成之前必须满足）。门控规则 DSL 支持四种操作符：

\begin{description}
\item[\texttt{var\_gte}] —— 命名变量必须大于或等于阈值。
\item[\texttt{arg\_in\_range}] —— 命名变量必须落在指定范围内。
\item[\texttt{flag\_true}] —— 布尔标志必须为 true。
\item[\texttt{time\_since}] —— 自指定事件以来必须经过最小时间。
\end{description}

% 证据：core/gating.py（evaluate_gate, evaluate_rule）

评估器包含显式的未知值处理：缺失、来源缺失或哨兵值变量会导致规则失败并附诊断原因，而非静默地视为 false。规则在首次失败时短路，失败索引和原因被传播到下游组件，包括确定性步骤提示和帮助循环上下文。
% 证据：core/gating.py, core/step_hint.py

门控规则以声明方式在程序包的配置中指定。F/A-18C 冷启动程序中的每个步骤都根据程序的操作要求被分配了前提门控和完成门控。分类为 BIOS 优先的步骤（S01、S03--S07、S09--S13、S21）主要依赖基于遥测的门控规则进行完成检测。分类为视觉优先的步骤（S08、S15、S18）以视觉事实证据补充门控评估，因为它们的完成条件涉及遥测中未反映的座舱显示状态。
% 证据：packs/fa18c_startup/pack.yaml（bios_priority_steps, vision_priority_steps）
% 声明：C-PROC-01, C-PROC-02

\subsection{增量去噪与聚合}

% 证据：adapters/delta_sanitizer.py, adapters/delta_aggregator.py, packs/fa18c_startup/delta_policy.yaml

原始遥测变量由于仿真器物理以高频率变化，即使没有执行任何学习者操作。为产生适合提示上下文的稳定状态表示，流水线应用两个连续的处理阶段。

\textbf{增量去噪。}
可配置的策略（在程序包的 \texttt{delta\_policy.yaml} 中指定）定义了逐变量阈值、防抖动窗口和过滤规则。由开关回弹、传感器噪声和 UDP 乱序传递引起的虚假变化被抑制。只有超过配置阈值并在防抖动窗口后持续存在的变化才被提升到去噪后的增量流中。
% 证据：adapters/delta_sanitizer.py, packs/fa18c_startup/delta_policy.yaml
% 声明：C-TELEM-01

\textbf{增量聚合。}
去噪后的增量在可配置时间窗口内累积，并被压缩为紧凑的、人类可读的摘要。该摘要作为 \texttt{recent\_deltas} 包含在帮助循环上下文中，为语言模型提供最近学习者操作和状态转换的简洁表示，而不暴露原始高频遥测数据。
% 证据：adapters/delta_aggregator.py

\subsection{确定性降级}

% 证据：core/step_hint.py, adapters/step_inference.py

当没有语言模型可用或模型输出验证失败时，程序运行时降级到确定性步骤推理。降级机制检查当前活跃步骤、其未满足的门控规则以及最近的 UI 目标，以产生限制在当前活跃步骤的文本提示。默认情况下，除非本地规则能够根据当前遥测状态证明一个唯一、有效的目标，否则在降级模式下不发送叠加高亮。这确保系统即使在没有模型的情况下也能保持运行和安全。
% 证据：core/step_hint.py, help_flow_en.md §5

% ===========================================================================
\section{知识溯源与来源策略}\label{sec:method-knowledge}

% 证据：adapters/knowledge_local.py, core/knowledge.py, knowledge_source_policy.yaml
% 声明：C-RAG-01

知识溯源子系统为教学运行时提供检索增强的程序性文档。它从已索引的文档集合中检索相关文本片段，并强制执行限制哪些文档可以呈现给模型的来源策略。

\subsection{BM25 检索}

% 证据：core/knowledge.py（BM25Retriever）, adapters/knowledge_local.py

知识检索使用 BM25，一种来自概率信息检索家族的词袋排序函数。文档使用 \texttt{index\_docs.py} 工具在句子级别进行索引，该工具将 F/A-18C NATOPS 飞行手册和补充冷启动文档预处理为本地 JSON 索引。
% 证据：tools/index_docs.py

在检索时，从当前程序状态构造一个溯源查询：活跃步骤标识符、其在步骤注册表中的描述以及最近的遥测观察。查询使用标准 BM25 加权函数以默认参数（$k_1 = 1.2$，$b = 0.75$）进行分词并与索引句子进行评分。返回 top-$k$ 得分片段（默认 $k = 5$），并作为可选上下文包含在帮助循环提示中。
% 证据：adapters/knowledge_local.py（build_grounding_query）

\subsection{来源策略}

% 证据：adapters/knowledge_source_policy.py, knowledge_source_policy.yaml

并非所有知识源都适用于每种教学上下文。声明式来源策略（在 \texttt{knowledge\_source\_policy.yaml} 中指定）定义了允许的文档来源允许列表。每个索引句子携带一个来源标识符，策略强制执行只有来自允许来源的片段才包含在检索结果中。
% 证据：adapters/knowledge_source_policy.py, knowledge_source_policy.yaml

来源策略服务于两个目的。首先，它防止模型暴露于与当前程序包相矛盾的程序变体——例如，当当前配置文件为陆基时的航母程序。其次，它将检索范围限制在经过筛选、已验证的文档中，降低因过时或不相关来源产生幻觉的风险。
% 证据：tests/test_knowledge_source_policy.py

知识子系统是可选的：如果不存在知识索引，帮助循环仅使用遥测和程序上下文继续。缺失的知识检索不得阻塞主教学流程。
% 证据：help_flow_en.md §3

% ===========================================================================
\section{视觉事实数据流水线}\label{sec:method-vlm-pipeline}

% 证据：tools/capture_vlm_dataset.py, tools/generate_vlm_prelabels.py,
%       tools/export_vision_sft_dataset.py, tools/train_qwen35_vlm_unsloth.py,
%       tools/benchmark_qwen35_vlm_facts.py
% 声明：C-VLM-01, C-VLM-03, C-VLM-07

视觉事实数据流水线将原始 DCS 座舱截图转换为结构化视觉事实提取模型。它包含七个阶段，涵盖数据捕获、人工标注、监督微调数据集生成、模型训练和基准测试评估。每个阶段产生带版本号、可检查的工件，实现从原始图像到基准测试指标的完全可追溯性。图\ref{fig:vlm-pipeline}提供了概览。

\begin{figure}[htbp]
  \centering
  % TODO: 插入 VLM 流水线图。
  % \includegraphics[width=\textwidth]{figures/vlm_pipeline.pdf}
  \caption{七阶段视觉事实数据流水线。捕获和预标注供给人工审查；经审查的标注驱动双语 SFT 导出；导出的数据训练 LoRA 适配器；基础模型和 LoRA 模型在独立留出集上进行基准测试。}
  \label{fig:vlm-pipeline}
\end{figure}

\subsection{第一阶段：DCS 截图捕获}

% 证据：tools/capture_vlm_dataset.py

第一阶段从 DCS World 捕获座舱截图。Lua 端 \texttt{SimTutor.lua} 脚本将三个主要多功能显示器——左 DDI、AMPCD 和右 DDI——渲染为单个复合面板图像。Python 端 \texttt{capture\_vlm\_dataset.py} 脚本编排捕获过程：它在可配置的时间间隔触发截图，将每次捕获与一个会话标识符和挂钟时间戳关联，并存储原始图像及副文件清单。
% 证据：tools/capture_vlm_dataset.py, DCS/Scripts/SimTutor/SimTutor.lua

选择复合面板格式是为了使训练、预标注、评估和运行时使用保持在同一输入分布上，避免训练-评估不匹配——即模型在训练期间接受三个分离区域裁剪，但在推理时接收单个复合图像。
% 证据：Doc/Vision/Reports/qwen35_vlm_finetune_report_EN.md §2

\subsection{第二阶段：初始 VLM 预标注}

% 证据：tools/generate_vlm_prelabels.py

每张捕获的复合面板图像由大型 VLM（通过 OpenAI 兼容 API 访问的 Qwen 397B 级模型）使用目标视觉事实本体作为提示契约进行预标注。预标注脚本将每张图像与一个结构化提示一起发送，该提示枚举本体中的所有事实、其定义以及三种可接受状态（\texttt{seen}、\texttt{not\_seen}、\texttt{uncertain}）。VLM 返回一个包含逐事实标签和证据注释的 JSON 对象。
% 证据：tools/generate_vlm_prelabels.py, tools/generate_vlm_prelabels_en.py, tools/generate_vlm_prelabels_zh.py

预标注作为种子步骤：它大规模生成初始标注，然后通过人工审查加以完善。预标注 VLM 是大型通用模型；它未在座舱领域进行微调，经常在模糊或稀有视觉状态上出错。下一阶段的人工审查纠正这些错误。

\subsection{第三阶段：Label Studio 人工审查}

% 证据：tools/generate_vlm_prelabels.py（Label Studio 导出路径）

预标注样本被导入 Label Studio，一个开源数据标注平台。每个任务向审查者呈现复合面板图像、VLM 对所有事实的预测标签以及一个自由文本摘要字段。审查者检查图像，纠正任何标注错误的事实，并为每处更正提供证据注释。审查界面强制要求每个事实接收三种可接受状态之一，并且摘要字段必须填写。
% 证据：Doc/Vision/Reports/qwen35_vlm_finetune_report_EN.md §3

人工审查是质量控制瓶颈：每个捕获会话包含 50 到 220 张图像，每张图像的审查时间取决于可见事实的数量和审查者对座舱显示的熟悉程度。所有经审查的样本都带有审计追踪，可追溯到原始图像、原始 VLM 预标注和审查者的更正。

\subsection{第四阶段：经审查的 JSONL 导出}

人工审查后，Label Studio 将更正后的标签导出为 JSONL 文件。每行包含任务标识符、原始图像路径、带逐事实状态和证据注释的经审查事实，以及经审查的摘要。此经审查的 JSONL 是所有下游阶段（SFT 导出、训练和基准测试）的权威真实标注工件。

\subsection{第五阶段：双语 SFT JSONL 生成}

% 证据：tools/export_vision_sft_dataset.py

经审查的 JSONL 由 \texttt{export\_vision\_sft\_dataset.py} 转换为 OpenAI 兼容的多模态对话 SFT 格式。每个经审查的样本生成两行训练数据——一行英文和一行中文——共享相同的复合面板图像和相同的事实标签，但在系统和用户指令文本上有所不同。
% 证据：tools/export_vision_sft_dataset.py, Doc/Vision/Reports/qwen35_vlm_finetune_report_EN.md §4

每个 SFT 行是一个三轮对话：一则系统消息指示模型仅输出 JSON，一则用户消息包含 base64 编码的图像和文本指令（枚举所有事实、其定义和决策边界），以及一则助手消息包含带有 \texttt{facts} 数组和 \texttt{summary} 字符串的真实标注 JSON 对象。助手输出中的每个事实条目恰好包含三个字段：\texttt{fact\_id}、\texttt{state} 和 \texttt{evidence\_note}。外部元数据字段如 \texttt{frame\_id}、\texttt{session\_id}、\texttt{confidence} 和图像路径被排除在训练目标之外，因为它们是数据集管理字段，而非视觉可推断属性。
% 证据：Doc/Vision/Reports/qwen35_vlm_finetune_report_EN.md §4.1--4.2

双语 SFT 增加的是指令多样性而非视觉多样性：英文和中文示例共享相同的图像和事实标签。其目的是提高两种语言下的 JSON 格式鲁棒性，并支持最终的双语运行时部署。
% 证据：Doc/Vision/Reports/qwen35_vlm_finetune_report_EN.md §4

\subsection{第六阶段：LoRA 微调}

% 证据：tools/train_qwen35_vlm_unsloth.py, models/*/train_summary.json

导出的 SFT JSONL 文件用于通过低秩适配（LoRA）微调基础 VLM。训练栈组合三个组件：Unsloth 用于 4 位 VLM 加载、LoRA 注入和视觉批次整理；PEFT 用于 LoRA 适配器定义；TRL \texttt{SFTTrainer} 用于监督微调循环。
% 证据：tools/train_qwen35_vlm_unsloth.py（FastVisionModel, get_peft_model, SFTTrainer）

训练脚本通过 Unsloth 的 \texttt{FastVisionModel.from\_pretrained} 以 4 位精度加载基础模型，对视觉和语言层（Qwen）或仅语言线性层（Gemma）应用 LoRA 适配器，并在双语 SFT 数据上进行训练。适配器权重与基础模型分开保存，便于高效存储和推理时加载。完整的训练配置细节见第\ref{sec:method-lora-config}节。
% 证据：models/qwen35_vlm_lora/README.md, models/gemma4_vlm_lora/.../train_summary.json

\subsection{第七阶段：基础模型与 LoRA 对比基准测试}

% 证据：tools/benchmark_qwen35_vlm_facts.py

最后阶段在留出数据集上评估微调后的 LoRA 适配器与未适配基础模型。基准测试脚本以 4 位推理模式加载每个模型（基础模型和基础模型加 LoRA），通过训练期间使用的相同提示发送每个留出样本，根据预期 JSON 模式解析和验证模型输出，并计算一套将模型预测与人工审查真实标注进行比较的度量。基准测试协议详见第\ref{sec:method-benchmark}节。
% 证据：tools/benchmark_qwen35_vlm_facts.py

\subsection{端到端可追溯性}

% 声明：C-VLM-03

该流水线的一个定义性方法论属性是端到端可追溯性。每个基准测试预测都可以通过训练数据、人工审查标签、VLM 预标注和原始 DCS 截图进行追溯。流水线在每个阶段产生带版本号的工件：捕获清单、预标注 JSONL 文件、Label Studio 导出文件、SFT JSONL 文件、LoRA 适配器权重、\texttt{train\_summary.json} 文件、基准测试预测以及基准测试度量摘要。这一监管链使得对任何观察到的模型行为进行严格审计成为可能，从单个事实误分类到聚合基准测试度量。
% 证据：benchmarks/qwen35_vlm_finetune/*/predictions_*.jsonl,
%       benchmarks/qwen35_vlm_finetune/*/errors_*.jsonl,
%       benchmarks/qwen35_vlm_finetune/*/benchmark_summary.json

% ===========================================================================
\section{数据集构建}\label{sec:method-datasets}

% 证据：datasets/*/stats.json
% 声明：C-DATA-01, C-VLM-07

\subsection{数据集概览}

视觉事实数据流水线在连续的捕获和审查周期中生成了六个数据集。表\ref{tab:datasets}总结了每个数据集的角色、样本数、事实本体版本、审查状态和导出语言。

\begin{table}[htbp]
  \centering
  \caption{数据集清单。所有样本数来自相应的 \texttt{stats.json} 文件。}
  \label{tab:datasets}
  \small
  \begin{tabular}{@{}lcccccc@{}}
  \toprule
  \textbf{数据集} & \textbf{角色} & \textbf{样本} & \textbf{本体} & \textbf{已审查} & \textbf{语言} \\
  \midrule
  Run-001 (\texttt{vision\_sft}) & 训练源（8-fact v1） & 180 & 8-fact v1 & 180 & en, zh \\
  \addlinespace
  Run-002 (\texttt{vision\_sft\_holdout\_run002}) & 留出（8-fact v1） & 50 & 8-fact v1 & 50 & en \\
  \addlinespace
  Run-002 newfacts (\texttt{vision\_sft\_holdout\_run002\_newfacts}) & 留出（13-fact v2） & 50 & 13-fact v2 & 50 & en \\
  \addlinespace
  Run-003 (\texttt{vision\_sft\_run003}) & 训练源（13-fact v2） & 220 & 13-fact v2 & 220 & en, zh \\
  \addlinespace
  Run-004 (\texttt{vision\_sft\_holdout\_run004\_random}) & 留出（13-fact v2） & 100 & 13-fact v2 & 100 & en \\
  \addlinespace
  Run-005 (\texttt{vision\_sft\_run005\_composition\_rebalance}) & 训练补充（13-fact v2） & 122 & 13-fact v2 & 122 & en, zh \\
  \bottomrule
  \end{tabular}
\end{table}

% 证据：datasets/vision_sft/stats.json（180 样本）
% 证据：datasets/vision_sft_holdout_run002/stats.json（50 样本）
% 证据：datasets/vision_sft_holdout_run002_newfacts/stats.json（50 样本）
% 证据：datasets/vision_sft_run003/stats.json（220 样本, missing_summary_count=82）
% 证据：datasets/vision_sft_holdout_run004_random/stats.json（100 样本）
% 证据：datasets/vision_sft_run005_composition_rebalance/stats.json（122 样本, missing_summary_count=25）

每个数据集均经过完全人工审查：每个样本都已在 Label Studio 中检查并更正。Run-003 和 Run-005 包含人工审查者将摘要字段留空的样本（分别有 82 和 25 个样本，如 \texttt{missing\_summary\_count} 字段所记录）。这些样本对训练和评估仍然有效，因为事实数组——主要监督目标——对每个样本中的每个事实都已完全标注。
% 证据：datasets/vision_sft_run003/stats.json（missing_summary_count=82）,
%       datasets/vision_sft_run005_composition_rebalance/stats.json（missing_summary_count=25）

\subsection{本体演进：从 8-Fact v1 到 13-Fact v2}

% 证据：packs/fa18c_startup/vision_facts.yaml, datasets/vision_sft/stats.json,
%       datasets/vision_sft_run003/stats.json
% 声明：C-VLM-07

初始的 8-fact v1 本体（用于 Run-001 和 Run-002）定义了八个事实：\texttt{fcs\_page\_visible}、\texttt{bit\_root\_page\_visible}、\texttt{bit\_page\_failure\_visible}、\texttt{right\_ddi\_fcsmc\_page\_visible}、\texttt{right\_ddi\_in\_test\_visible}、\texttt{fcs\_bit\_result\_visible}、\texttt{ins\_alignment\_page\_visible} 和 \texttt{ins\_go}。
% 证据：datasets/vision_sft/stats.json（fact_state_counts 键）

v1 本体在初始基准分析中暴露出两个结构性弱点。首先，它混合了抽象层次：一些事实描述低层视觉观察（如 \texttt{fcs\_page\_visible}），而另一些则编码了更高层的程序性判断（如 \texttt{ins\_go}，它代表 INS 校准过程的完成信号而非直接可观察的视觉特征）。其次，若干事实共享重叠语义：\texttt{bit\_root\_page\_visible} 和 \texttt{bit\_page\_failure\_visible} 可能在某些座舱状态下重叠，而 \texttt{right\_ddi\_fcsmc\_page\_visible} 将区域标识符与事实语义混淆。
% 证据：Doc/Vision/Reports/qwen35_vlm_finetune_report_EN.md §8.2

13-fact v2 本体（用于 Run-003、Run-004、Run-005 及当前运行时）重新设计并扩展了事实集。v2 事实为：\texttt{tac\_page\_visible}、\texttt{supt\_page\_visible}、\texttt{fcs\_page\_visible}、\texttt{fcs\_page\_x\_marks\_visible}、\texttt{bit\_root\_page\_visible}、\texttt{fcsmc\_page\_visible}、\texttt{fcsmc\_intermediate\_result\_visible}、\texttt{fcsmc\_in\_test\_visible}、\texttt{fcsmc\_final\_go\_result\_visible}、\texttt{hsi\_page\_visible}、\texttt{hsi\_map\_layer\_visible}、\texttt{ins\_grnd\_alignment\_text\_visible} 和 \texttt{ins\_ok\_text\_visible}。
% 证据：packs/fa18c_startup/vision_facts.yaml, tools/benchmark_qwen35_vlm_facts.py（CORE_FACT_IDS）

v2 重新设计通过三种策略解决了 v1 的弱点。首先，模糊的 \texttt{ins\_go} 事实被分解为两个低层、直接可观察的事实：\texttt{ins\_grnd\_alignment\_text\_visible}（要求字面 GRND 校准块文本）和 \texttt{ins\_ok\_text\_visible}（要求校准块附近清晰可读的 OK 文本）。其次，区域标识符从事实名称中移除（\texttt{right\_ddi\_fcsmc\_page\_visible} 变为 \texttt{fcsmc\_page\_visible}），并为 TAC 和 SUPT 菜单添加了页面可见性事实以覆盖左 DDI。第三，FCS-MC 状态序列被分解为更细粒度的四状态链（\texttt{page}、\texttt{intermediate\_result}、\texttt{in\_test}、\texttt{final\_go\_result}），替换了单一的 \texttt{fcs\_bit\_result\_visible} 事实。第四，HSI 相关事实被拆分为 \texttt{hsi\_page\_visible} 和 \texttt{hsi\_map\_layer\_visible}，FCS X 标记事实与 FCS 页面可见性分离。
% 证据：Doc/Vision/Reports/qwen35_vlm_finetune_report_EN.md §9,
%       Doc/Vision/Reports/vlm_backbone_comparison_summary_EN.md

\subsection{训练方案：Run-003 + Run-005x2}

% 证据：models/qwen35_vlm_lora/full_qwen35_9b_base_bilingual_run003_plus_run005x2_v1/train_summary.json

13-fact v2 模型的主要训练方案将 Run-003（220 个经审查样本，双语导出生成 220 条英文和 220 条中文训练行）与包含两次的 Run-005 组合再平衡数据（122 个经审查样本，双语导出生成 $122 \times 2 = 244$ 条英文和 $122 \times 2 = 244$ 条中文训练行）相结合。得到的训练集包含 928 行。
% 证据：models/qwen35_vlm_lora/full_qwen35_9b_base_bilingual_run003_plus_run005x2_v1/train_summary.json（train_rows=928）

Run-005 专门作为组合再平衡数据集采集。对 Run-003 事实状态分布的分析揭示了严重的类别不平衡：若干关键事实（如 \texttt{fcsmc\_intermediate\_result\_visible}、\texttt{fcsmc\_in\_test\_visible}、\texttt{fcsmc\_final\_go\_result\_visible}、\texttt{ins\_ok\_text\_visible}）在训练数据中具有极少的 \texttt{seen} 示例。Run-005 设计了目标座舱状态转换以过采样这些代表性不足的视觉状态。
% 证据：datasets/vision_sft_run003/stats.json（fact_state_counts 显示 ins_ok_text_visible 的 seen 计数低至 12）,
%       datasets/vision_sft_run005_composition_rebalance/stats.json（fact_state_counts 显示改善的平衡性）

Run-005 的双重计数（Run-005x2）进一步放大了这些稀有状态在训练分布中的代表性。该训练运行的输入文件为：来自 Run-003 的 \texttt{sft\_en.jsonl} 和 \texttt{sft\_zh.jsonl}，加上来自 Run-005 的 \texttt{sft\_en.jsonl} 和 \texttt{sft\_zh.jsonl} 各两份副本，共计六个输入文件。
% 证据：models/qwen35_vlm_lora/full_qwen35_9b_base_bilingual_run003_plus_run005x2_v1/train_summary.json（input_files 数组）

未使用内部评估划分：训练配置设置 \texttt{eval\_rows = 0}。所有评估均在独立留出数据集上执行。
% 证据：models/qwen35_vlm_lora/full_qwen35_9b_base_bilingual_run003_plus_run005x2_v1/train_summary.json（eval_rows=0）

\subsection{留出集独立性验证}

% 证据：Doc/Vision/Reports/vlm_backbone_comparison_summary_EN.md
% 声明：C-VLM-03a

留出集独立性是一项关键的方法论要求：基准测试必须测量跨会话泛化，而非分布内记忆。构建了两个留出集：

\begin{itemize}
\item \textbf{Run-002}（及其 13-fact 重标注 \texttt{run002\_newfacts}）：一个与所有训练会话不同的新 DCS 捕获会话。包含 50 个经审查样本。
\item \textbf{Run-004 random}：一个进一步独立的捕获会话，采用均匀随机采样的座舱状态。包含 100 个经审查样本。
\end{itemize}

% 证据：Doc/Vision/Reports/vlm_backbone_comparison_summary_EN.md

每个留出集与训练数据之间的完全重叠经验证为零（0）。重叠检查比较了训练和留出划分中的图像工件路径和样本标识符。没有留出样本出现在任何训练输入文件中。
% 证据：Doc/Vision/Reports/vlm_backbone_comparison_summary_EN.md（完全重叠 = 0）

相比之下，早期 8-fact v1 实验中使用的 Run-001 基准测试是一个\textbf{受污染开发集}：训练示例所来源的同一数据池被用于评估。Run-001 结果仅为回归分析报告，并与独立留出结果明确区分。
% 证据：benchmarks/qwen35_vlm_finetune/base_vs_lora_current180_v1/benchmark_summary.json（benchmark_kind=contaminated_dev_set）

% ===========================================================================
\section{LoRA 微调配置}\label{sec:method-lora-config}

% 证据：models/*/train_summary.json, tools/train_qwen35_vlm_unsloth.py,
%       tools/train_gemma4_vlm_unsloth.py
% 声明：C-TRAIN-01

\subsection{基础模型}

两种基础 VLM 架构在相同训练数据方案下进行了微调以实现骨干对比：

\begin{description}
\item[Qwen/Qwen3.5-9B-Base] —— 一个 90 亿参数的视觉语言模型，支持多模态输入（图像 + 文本）和文本生成。这是主要的实验骨干。
% 证据：tools/train_qwen35_vlm_unsloth.py（model_name 默认值）

\item[google/gemma-4-31B] —— 来自 Google DeepMind 的 310 亿参数视觉语言模型。这是补充对比骨干，用于检验相同训练数据方案是否可跨模型家族泛化。
% 证据：models/gemma4_vlm_lora/full_gemma4_31b_base_bilingual_run003_plus_run005x2_v1/train_summary.json
\end{description}

两个模型均通过 Unsloth 的 \texttt{FastVisionModel.from\_pretrained} 以 4 位精度加载。对 Qwen，LoRA 适配器应用于视觉和语言层（\texttt{finetune\_vision\_layers=True}、\texttt{finetune\_language\_layers=True}）。对 Gemma，LoRA 被限制在所有线性模块（\texttt{lora\_target\_modules=all-linear}），且视觉层未微调（\texttt{finetune\_vision\_layers=False}）。
% 证据：tools/train_qwen35_vlm_unsloth.py（get_peft_model 调用）,
%       models/gemma4_vlm_lora/full_gemma4_31b_base_bilingual_run003_plus_run005x2_v1/train_summary.json

\subsection{训练栈}

% 证据：tools/train_qwen35_vlm_unsloth.py

训练栈包含三个集成组件：

\begin{enumerate}
\item \textbf{Unsloth} —— 处理 4 位模型加载、LoRA 注入、通过 \texttt{UnslothVisionDataCollator} 进行视觉批次整理，以及 FP16/BF16 自动精度选择。
\item \textbf{PEFT（参数高效微调）} —— 定义 LoRA 适配器格式，包括秩、alpha、dropout 和目标模块选择。
\item \textbf{TRL SFTTrainer} —— 以可配置的批次大小、梯度累积、轮数、学习率、预热和权重衰减运行监督微调循环。
\end{enumerate}

% 证据：tools/train_qwen35_vlm_unsloth.py（导入和 trainer_args）

\subsection{超参数配置}

表\ref{tab:lora-hyperparams}列出了主要训练运行（Run-003 + Run-005x2）使用的精确超参数，在两个骨干间匹配。所有值取自 \texttt{train\_summary.json} 工件。

\begin{table}[htbp]
  \centering
  \caption{LoRA 微调超参数，在 Qwen3.5-9B 和 Gemma-4-31B 间匹配。}
  \label{tab:lora-hyperparams}
  \small
  \begin{tabular}{@{}lll@{}}
  \toprule
  \textbf{参数} & \textbf{Qwen3.5-9B} & \textbf{Gemma-4-31B} \\
  \midrule
  \texttt{max\_seq\_length} & 4096 & 4096 \\
  \texttt{num\_train\_epochs} & 4.0 & 4.0 \\
  \texttt{learning\_rate} & 2$\times 10^{-4}$ & 2$\times 10^{-4}$ \\
  \texttt{per\_device\_train\_batch\_size} & 1 & 1 \\
  \texttt{gradient\_accumulation\_steps} & 4 & 4 \\
  \texttt{lora\_r} & 16 & 16 \\
  \texttt{lora\_alpha} & 16 & 16 \\
  \texttt{lora\_dropout} & 0.0 & 0.0 \\
  \texttt{seed} & 3407 & 3407 \\
  \texttt{load\_in\_4bit} & True & True \\
  \texttt{warmup\_ratio} & 0.1 & 0.1 \\
  \texttt{weight\_decay} & 0.01 & 0.01 \\
  \texttt{train\_rows} & 928 & 928 \\
  \texttt{eval\_rows} & 0 & 0 \\
  \midrule
  \textbf{骨干特定} & & \\
  \midrule
  视觉 LoRA & 启用 & 禁用 \\
  LoRA 目标模块 & vision + language + attention + MLP & all-linear \\
  对话模板 & Qwen 默认 & gemma-4 \\
  GPU 内存利用率 & 0.6 & 0.95 \\
  \midrule
  \textbf{训练结果} & & \\
  \midrule
  训练运行时间（秒） & 6419.16 & 8049.65 \\
  训练损失（最终） & 0.2156 & 0.0474 \\
  \bottomrule
  \end{tabular}
\end{table}

% 证据：models/qwen35_vlm_lora/full_qwen35_9b_base_bilingual_run003_plus_run005x2_v1/train_summary.json
% 证据：models/gemma4_vlm_lora/full_gemma4_31b_base_bilingual_run003_plus_run005x2_v1/train_summary.json
% 证据：Doc/Vision/Reports/vlm_backbone_comparison_summary_EN.md

最终训练损失值（Qwen 为 0.2156，Gemma 为 0.0474）不应在骨干间直接比较，因为分词器行为、对话模板结构、目标序列分布和模型参数化各不相同。4 个轮数的选择是作为小数据领域适配的折衷：仅有 928 个训练样本，更少的轮数可能无法可靠地教会模型固定的 JSON 模式和座舱事实边界，而更多的轮数则有过度拟合训练会话的风险。
% 证据：Doc/Vision/Reports/qwen35_vlm_finetune_report_EN.md §5

LoRA 秩 $r = 16$ 和 alpha $\alpha = 16$ 提供适度的适配能力：足以学习座舱布局模式和 JSON 格式，但不像高容量适配器（$r \geq 64$）那样可能更激进地过度拟合小数据集。在主要运行中将 Dropout 设置为零，在首次适配实验中优先考虑可学习性而非正则化。
% 证据：Doc/Vision/Reports/qwen35_vlm_finetune_report_EN.md §5

\subsection{双语 SFT 策略}

% 证据：tools/export_vision_sft_dataset.py

双语 SFT 策略为每个经审查样本生成两行训练数据：一行使用英文系统和用户指令，一行使用中文指令。两行共享相同的图像（base64 编码）和相同的助手目标（真实标注事实 JSON）。英文助手目标被复用于中文行；模型的任务是从图像中提取视觉事实，而非生成语言特定的自然语言描述。
% 证据：tools/export_vision_sft_dataset.py, Doc/Vision/Reports/qwen35_vlm_finetune_report_EN.md §4

这种方法在不增加额外人工标注工作量的情况下提高了跨语言指令遵循的鲁棒性。然而，它不能替代额外的座舱状态截图或困难负样本视觉数据。
% 证据：Doc/Vision/Reports/qwen35_vlm_finetune_report_EN.md §8.3

\subsection{组合再平衡策略}

% 证据：datasets/vision_sft_run005_composition_rebalance/stats.json

组合再平衡策略（Run-005）解决了主要训练源（Run-003）中观察到的类别不平衡。在 Run-003 中，若干事实具有极低的正例（\texttt{seen}）计数：\texttt{ins\_ok\_text\_visible} 在 220 个样本中仅有 12 个 \texttt{seen} 示例（5.5\%）；\texttt{fcs\_page\_x\_marks\_visible} 有 16 个（7.3\%）；\texttt{fcsmc\_intermediate\_result\_visible}、\texttt{fcsmc\_in\_test\_visible} 和 \texttt{fcsmc\_final\_go\_result\_visible} 各有 18 个（8.2\%）。
% 证据：datasets/vision_sft_run003/stats.json（fact_state_counts）

Run-005 捕获了 122 个额外样本，针对这些稀有事实为 \texttt{seen} 的座舱状态。在 Run-005 增强后，组合训练集提供了对先前代表性不足事实的改进覆盖：\texttt{supt\_page\_visible}（组合 seen 从 20 到 54）、\texttt{fcsmc\_page\_visible}（从 55 到 111）、\texttt{hsi\_page\_visible}（从 106 到 215）以及 \texttt{ins\_grnd\_alignment\_text\_visible}（从 58 到 130）。
% 证据：datasets/vision_sft_run003/stats.json, datasets/vision_sft_run005_composition_rebalance/stats.json

Run-005 数据在训练方案中被包含两次（Run-005x2），进一步放大了这些代表性不足视觉状态的信号。
% 证据：models/qwen35_vlm_lora/full_qwen35_9b_base_bilingual_run003_plus_run005x2_v1/train_summary.json

% ===========================================================================
\section{基准测试协议}\label{sec:method-benchmark}

% 证据：tools/benchmark_qwen35_vlm_facts.py
% 声明：C-VLM-04, C-VLM-05, C-VLM-06

\subsection{度量指标}

基准测试使用以下度量比较模型预测与人工审查真实标注标签：

\begin{description}
\item[\texttt{json\_valid\_rate}] —— 模型输出可被解析为 JSON 对象的样本比例。此处的失败表明模型产生了非 JSON 文本（如 markdown、自然语言评论或格式错误的括号）。
\item[\texttt{schema\_valid\_rate}] —— 解析后的 JSON 符合预期模式的样本比例：顶层对象仅包含 \texttt{summary} 和 \texttt{facts} 键，\texttt{facts} 数组中每个条目恰好包含 \texttt{fact\_id}、\texttt{state} 和 \texttt{evidence\_note}，所有 13 个事实 ID 均恰好出现一次，所有状态均在 \texttt{\{seen, not\_seen, uncertain\}} 中。
\item[\texttt{fact\_accuracy}] —— 所有样本中所有事实的三类准确率，计算为与真实标注标签匹配的单个事实预测的比例。
\item[\texttt{macro\_f1}] —— 跨三种状态（\texttt{seen}、\texttt{not\_seen}、\texttt{uncertain}）的宏平均 F1 分数，逐事实计算后跨所有 13 个事实平均。
\item[\texttt{seen\_f1}] —— 正类 \texttt{seen} 的 F1 分数，逐事实计算后跨所有 13 个事实平均。这是一个特别重要的指标，因为遗漏 \texttt{seen} 状态（假阴性）或幻觉出 \texttt{seen} 状态（假阳性）直接影响下游程序状态推理。
\item[\texttt{sample\_exact\_match}] —— 所有 13 个事实与真实标注标签完全匹配的样本比例。单个事实的误分类计为样本级失败。
\item[\texttt{critical\_false\_positive\_count}] —— 分类为\emph{关键}的事实子集（6 个事实的子集：\texttt{fcs\_page\_x\_marks\_visible}、\texttt{fcsmc\_intermediate\_result\_visible}、\texttt{fcsmc\_in\_test\_visible}、\texttt{fcsmc\_final\_go\_result\_visible}、\texttt{ins\_grnd\_alignment\_text\_visible}、\texttt{ins\_ok\_text\_visible}）中，当真实标注为 \texttt{not\_seen} 或 \texttt{uncertain} 时被预测为 \texttt{seen} 的错误计数。
\end{description}

% 证据：tools/benchmark_qwen35_vlm_facts.py（compute_metrics, CRITICAL_FACT_IDS）

关键假阳性被单独跟踪，因为它们代表了教学上下文中最具后果性的错误类型：错误地报告一个完成关键状态为已观察而实际上并非如此。例如，一个错误地认为 INS 校准已完成（基于 \texttt{ins\_ok\_text\_visible} 假阳性）的导师可能跳过必要的指导，潜在地导致程序性错误。
% 证据：Doc/Vision/Reports/qwen35_vlm_finetune_report_EN.md §6

基准测试还计算逐事实混淆矩阵以及逐事实的精确率、召回率和 F1 分数，这些存储在 \texttt{fact\_scores.csv} 中，并附带自动生成的图表（总体准确率、样本完全匹配、逐事实宏 F1、逐事实 seen F1、关键假阳性以及逐事实混淆矩阵）。
% 证据：tools/benchmark_qwen35_vlm_facts.py（_make_charts, _write_fact_scores_csv）

\subsection{基础模型与 LoRA 对比方法论}

% 证据：tools/benchmark_qwen35_vlm_facts.py

基准测试协议比较两种模型配置在相同真实标注标签上的表现：

\begin{enumerate}
\item \textbf{基础模型}：未适配的基础 VLM，以 4 位推理模式加载，不带任何 LoRA 适配器。
\item \textbf{LoRA 模型}：相同基础 VLM，以 4 位推理模式加载，通过 PEFT 的 \texttt{PeftModel.from\_pretrained} 加载 LoRA 适配器权重。
\end{enumerate}

两种配置接收相同的提示、相同的图像和相同的生成参数：temperature = 0.0（贪婪解码）、\texttt{max\_new\_tokens} = 1024、seed = 3407。基准测试独立运行每个样本，将预测序列化到 JSONL 中，并在没有任何事后过滤或排除的情况下计算度量。
% 证据：tools/benchmark_qwen35_vlm_facts.py（_run_predictions, temperature=0.0）

对比脚本（\texttt{comparison.json}）计算增量度量（LoRA 减去基础模型）和一个自动化推荐标志（\texttt{recommend\_simtutor\_chain\_validation}），该标志在 LoRA 适配器同时提升 JSON 有效性、事实准确率、seen F1，且不增加关键假阳性时激活。
% 证据：tools/benchmark_qwen35_vlm_facts.py（_compare_metrics）

\subsection{留出定义与基准测试类别}

% 证据：benchmarks/qwen35_vlm_finetune/*/benchmark_summary.json

每个基准测试运行被赋予一个 \texttt{benchmark\_kind}，决定其结果应如何解读：

\begin{description}
\item[\texttt{contaminated\_dev\_set}] —— 评估集从与训练集相同的数据池中抽取。结果表明模型是否学会了模式和分布内视觉边界，但不能支持泛化声明。用于 Run-001 基准测试。
% 证据：benchmarks/qwen35_vlm_finetune/base_vs_lora_current180_v1/benchmark_summary.json

\item[\texttt{heldout\_new\_session}] —— 评估集是一个完全独立的捕获会话，与任何训练数据零重叠。结果支持跨会话泛化声明。这是主要评估类别。Run-002 newfacts 和 Run-004 random 基准测试属于此类别。
% 证据：benchmarks/qwen35_vlm_finetune/run003_plus_run005x2_vs_base_holdout_run002_newfacts_v1/benchmark_summary.json（benchmark_kind=holdout_run002_newfacts）
\end{description}

受污染与留出基准测试之间的区分在整个基准测试基础设施中被严格执行，并反映在结果章节（第\ref{ch:evaluation}章）中。
% 证据：paper_outline.md §4

\subsection{基准测试工件}

% 证据：benchmarks/qwen35_vlm_finetune/run003_plus_run005x2_vs_base_holdout_run002_newfacts_v1/

每个基准测试运行生成一个标准化的工件目录：

\begin{itemize}
\item \texttt{benchmark\_summary.json} —— 顶层摘要，包含对比增量和图表路径。
\item \texttt{metrics\_base.json} / \texttt{metrics\_lora.json} —— 完整度量字典，包括逐事实分数和混淆矩阵。
\item \texttt{predictions\_base.jsonl} / \texttt{predictions\_lora.jsonl} —— 逐样本模型输出，每行一个 JSON 对象。
\item \texttt{errors\_base.jsonl} / \texttt{errors\_lora.jsonl} —— 单个预测错误，包含金标签、预测标签和原始模型文本。
\item \texttt{comparison.json} —— 增量度量和推荐标志。
\item \texttt{fact\_scores.csv} —— 逐事实逐模型度量的 CSV 表。
\item \texttt{report.md} —— 人类可读的 markdown 摘要。
\item \texttt{charts/} —— 自动生成图表：总体准确率、样本完全匹配、逐事实宏 F1、逐事实 seen F1、关键假阳性，以及基础模型和 LoRA 模型的逐事实混淆矩阵。
\end{itemize}

% 证据：benchmarks/qwen35_vlm_finetune/run003_plus_run005x2_vs_base_holdout_run002_newfacts_v1/benchmark_summary.json

这些工件构成了第\ref{ch:evaluation}章中报告的技术结果的完整证据链。通过使用相同的经审查 JSONL 和 LoRA 适配器权重重新运行基准测试脚本，每个度量、图表和错误记录都是可复现的。

% ===========================================================================
\section{本章小结}\label{sec:method-summary}

本章描述了 SimTutor 系统的四个方法论支柱：基于遥测证据的程序运行时、以 BM25 检索和来源策略为基础的知识溯源、从截图捕获到基准测试的七阶段视觉事实数据流水线，以及在两个骨干架构间匹配超参数的 LoRA 微调协议。该方法论的特征在于端到端可追溯性——每个基准测试度量可通过训练数据、人工审查标签和原始截图进行追溯——以及受污染开发基准测试与独立留出基准测试之间的严格分离。这些方法论选择直接支撑了第\ref{ch:evaluation}和第\ref{ch:discussion}章中呈现的声明和结果。
